\documentclass[10pt]{article}
\usepackage[letterpaper,margin=.38in]{geometry}
\usepackage{amsmath,amssymb,multicol}
\pdfinfoomitdate=1\pdftrailerid{}
\pagestyle{empty}\setlength{\parindent}{0pt}\setlength{\parskip}{2pt}
\setlength{\columnsep}{13pt}\setlength{\jot}{0pt}
\setlength{\abovedisplayskip}{3pt}\setlength{\belowdisplayskip}{3pt}
\let\bigsum\sum\renewcommand{\sum}{\bigsum\nolimits}
\newcommand{\fl}[1]{\left\lfloor#1\right\rfloor}
\newcommand{\rg}[2]{\{#1,\ldots,#2\}}
\begin{document}
\begin{center}
{\large\bfseries A complete fifteen-variable consistency equation}\\[-1pt]
{\small Strictly positive integer unknowns; exact degree 157; a 66-operation evaluator.}
\end{center}\vspace{-8pt}
\begin{multicols}{2}
\fontsize{9}{9.5}\selectfont
The only unknowns are
\[
J,F,\alpha,z,f,h,i,T,s,w,\tau,\eta,\zeta,y,\gamma\in\mathbb Z_{>0}.
\]
All other symbols in the equation are polynomial abbreviations:
\[
\begin{gathered}
q_-=\beta J,\quad q=q_-+1,\quad C=q-F-\alpha,\\
X=wq,\quad Y=sq^3,\quad E=XY,\quad k=\eta+\zeta,\\
\ell=EkY,\quad c=kY+\eta,\quad a=E+Y,\\
H=4a+3,\quad\Delta=a^2+H,\quad D=X+ac+\gamma H,\\
R=[q_-(q-F)+\alpha](q^2-1)+(M_C+qM_F)J,\\
V=f^2(cT-R)-c,\quad S=i\Delta c^2,\\
N_0=\tau^2-\ell(\ell+k),\quad N_1=D^2-\Delta c^2,\\
N_a=S^2(V^2-y^2)+y^2,\quad N_k=k-hE-R,\\
N_t=(K+w)C+(q-F)-zq_-,\quad N_s=\Delta f^2-S^2.
\end{gathered}
\]
The single equation is
\[
\boxed{N_0N_1N_aN_kN_tN_s-\Delta=0.}
\]

\textbf{The four fixed integers.}
Every index, set, sum, product and floor below defines a constant, not
another unknown. Put
\[
\rho=2^{67},\quad B=\rho^{8807},\quad\beta=B-1,\quad
G_m(x)=\sum_{j=0}^{m-1}x^j.
\]
Define
\[
\begin{split}
\mathcal E={}&\rg0{942}\cup\rg{983}{992}\cup\rg{1003}{1011}\\
&\cup\rg{1013}{1031}\cup\{1063,1064,1095,2195,3296\}.
\end{split}
\]
Then
\begin{align*}
M_C={}&B-1-\sum_{e\in\mathcal E\setminus\{1\}}\rho^e-2\rho^{1063},\\
M_F={}&B-1+4+(2^{51}-2)(\rho^{4393}+\rho^{5495})\\
&+\rho^{1013}G_{10}(\rho)+\rho^{1065}G_9(\rho)
  +\rho^{2195}+\rho^{3296}\\
&+(2^{65}-1)\rho^{8792}G_{13}(\rho)+(2^{19}-1)\rho^{8805},\\
K={}&\rho^{30}+\rho^{62}+\rho^{1100}+\rho^{2201}+\rho^{8849}\\
&+(\rho^{4393}+\rho^{5495})\Phi(1/\rho),
\end{align*}
where
\[
\begin{split}
\Phi(x)={}&(1+2^{51})G_{943}(x)
 -2^{10}\sum_{i=0}^{942}v(t_i)x^i\\
&+\sum_{(b,e,m)\in\mathcal C}2^b x^e G_m(2x),\\
\mathcal C={}&\{(10,1023,8),(18,1013,9),(27,1003,8),\\
&(35,983,9),(44,1031,1),(45,1022,1),(46,1011,1),\\
&(47,992,1),(48,1095,1),(49,2195,1),(50,3296,1)\}.
\end{split}
\]
The rational intermediate in K cancels exactly; all four results are positive integers.

\textbf{The 943 indexed fixed tiles.}
Use the actions $u_j,v_j,d_j$ decoded below. Let
\[
\begin{gathered}
h(q,b)=2(q+1)+b,\quad
\mathcal R=\{u_j:d_j=1,\ 0\le j<430\},\\
\mathcal L=\{u_j:d_j=0,\ 0\le j<430\},\\
r_q=6+|\mathcal R\cap\rg0{q-1}|,\quad
l_q=205+|\mathcal L\cap\rg0{q-1}|.
\end{gathered}
\]
An interval ending at $-1$ is empty. The table defines $t_i$ for
$0\le i\le942$; all index ranges are fixed:
\[
\setlength{\arraycolsep}{1.4pt}\begin{array}{c|l}
\text{index }i&\text{tile }t_i=(w,n,e,s)\\\hline
0,1&(2,435,3,1),\ (5,0,0,0)\\
j+2\ (j=2q+b<430)&((1-d_j)l_{u_j},h(q,b),d_jr_{u_j},v_j)\\
432+b\ (b=0,1)&(0,h(215,b),0,h(215,b))\\
434+b\ (b=0,1)&(0,b,0,b)\\
424+2r_q+b\ (q\in\mathcal R,b=0,1)&(r_q,b,0,h(q,b))\\
424+2l_q+b\ (q\in\mathcal L,b=0,1)&(0,b,l_q,h(q,b))\\
924+j\ (0\le j<4)&(0,j+430\fl{j/2},0,438)\\
928&(2,434,2,0)\\
929&(3,436,4,4)\\
930&(4,436,4,0)\\
497+t\ (t=434,435,436)&(1,438,1,t)\\
\substack{934+3a+b-6\fl{a/4}\\a\in\{0,1,4\},\ b=0,1,2}&(a,437,b,437)
\end{array}
\]
For $t=(w,n,e,s)$ put
\[
v(t)=w+2^8n+2^{17}e+2^{25}s+127\,2^{34}.
\]

\textbf{A closed decoder for the 430 actions.}
Put $D_b(A,k)=\fl{A/b^k}\bmod b$, where
$a\bmod b=a-b\fl{a/b}$. For $0\le q<215$, $b=0,1$, $j=2q+b$:
\[
\begin{gathered}
N=\mathcal A\bmod2^{636},\quad
T_0=\fl{\mathcal A/2^{636}}\bmod2^{1675},\\
E_0=\fl{\mathcal A/2^{2311}},\\
\nu_j=D_2(N,j),\quad \sigma_j=\sum_{k=0}^{j-1}\nu_k,\\
m_q=D_2(N,430+q),\quad \tau_q=\sum_{k=0}^{q-1}m_k,\\
F_q=14(1-m_q)+m_qD_{16}(E_0,\tau_q),\\
u_j=\nu_j(1+\sigma_j)+(1-\nu_j)D_{216}(T_0,j-\sigma_j),\\
v_j=D_2(F_q,2b),\quad d_j=D_2(F_q,2b+1).
\end{gathered}
\]
Empty sums are zero. Here $(v_0,d_0,u_0)=(1,1,1)$; the constructor tile
uses this actual first target, including its state renaming.

\textbf{The integer $\mathcal A$.}
The alphabet, in order of digit values $0,\ldots,61$, is
{\fontsize{7.7}{8.2}\selectfont\ttfamily 0123456789ABCDEFGHIJKLMNOPQRSTUVWXYZabcdefghijklmnopqrstuvwxyz}.\par
Concatenate the following 436 characters as digits $e_0,\ldots,e_{435}$;
put $\mathcal A=\sum_{j=0}^{435}e_j62^{435-j}$.\par
\begingroup\fontsize{7.7}{8.2}\selectfont\ttfamily
BxEnhvy6NipzkWmROo0GsA3uK3twWLxscLq20UTumy5Cdu461LwD6g7RDNzF8Y\par
iYDFTpB8ehAypd7d6OsOL4APamnJeNXQypUezCmW9n4oXP6XfgmiWuAWZ8ZJ7U\par
OP7gu4kVdlGZqZHULQMar1dXNeJW4i7oNVI8usp8BlovjjZbwPCovZI328fRUq\par
4q92R81s5GjWQaQfGhkq9UYMK3rRMaBmsOIBkJqc2cj2VRIZ6mjzkBIOkW9SjZ\par
FzeA6J5u4j3sBL3SdgwXzM9mHvQivu8EJ1TBTpRhHDqZIbmYgwQ8VqTAFotl6F\par
We1YC554GKrVHsA8oUGKSA6F6k5KREBSD0Tgd7bFALY84LLbpw574BQyvdoHrl\par
0Oe35uwDf15DpT6VJWZBjSbHPpPAINYTte0eUxnTFTKBlGj8Am6sb1gLldh8WZ\par
8P\par
\endgroup
\normalfont\fontsize{7.7}{8.2}\selectfont
The ordinary state/color-conjugacy, native and Pell arguments give
unsolvability iff $\mathrm{Con}(\mathrm{ZFC})$.
Native calibration and internal-ZFC formalization remain open in Lean.
This page is a complete finite specification, not an unsolvability proof.
\end{multicols}
\end{document}
